\section{Discussion}
\label{sec:discussion}

\pending{Interpretation of the results: what dominates the remaining error, and
what it means for users of Hinglish meetings.}

\subsection{Limitations}
\paragraph{Script-based language identification}
The language tagger reads script, not language. Hindi that the recognizer
writes in Latin script is tagged \textsc{English}, so its accuracy is bounded by
how consistently the recognizer uses Devanagari for Hindi.
\pending{Quantify from the confusion matrix.}

\paragraph{Segment-level timing}
Speakers are assigned per recognized segment, so a segment spanning a change of
speaker is attributed entirely to one of them. Word-level timestamps would
remove this error at some cost in speed.

\paragraph{Speaker resolution from short or mixed windows}
Embeddings of windows shorter than 1\,s are not trusted and are left unresolved,
which increases the number of unresolved turns in rapid exchanges. Resolving a
whole chunk with one embedding is faster but blends the voices of different
speakers in the same chunk; Table~\ref{tab:ablation_results} measures this
trade-off.

\paragraph{Dataset size and conditions}
\pending{N} clips recorded by \pending{participants} with \pending{devices} do
not cover the range of accents, room acoustics and meeting styles found in
practice, and the confidence intervals reflect this.

\paragraph{Language-model variability}
Language-model outputs vary between runs even at low temperature; we report the
spread over repeated runs rather than a single run.

\subsection{Threats to validity}
\pending{Annotator bias (annotators are also authors?); reference script
conventions; the choice of the 0.5 Jaccard threshold for action items;
generalization beyond Hindi--English.}
